\chapter{Outils connectés dans le dépôt \texttt{paas}}
\label{ch:outils}

\section{Introduction}

Les paragraphes qui suivent ne sont~\textbf{pas} un catalogue généraliste de l'éditeur~: ils reprennent les briques~\emph{explicitement nominées dans \texttt{paas/README.md}, \texttt{paas/docs/STACK\_OVERVIEW.md} et reliées depuis le fichier \texttt{.env.example} du frontal}, car ce sont-elles qui servent lorsque nous branchons~(ou projetons)~l'infrastructure hébergeant~\texttt{paas/frontend}.

\section{Cartographie par rapport aux variables d'environnement}

Une ligne typique~(voir fichier schématisé en annexe)~suit~\texttt{HARBOR\_*},~\texttt{JENKINS\_*} ou~\texttt{TEKTON\_*},~\texttt{GITOPS\_*},~\texttt{ARGOCD\_*}. Les autres outillage~(SonarQube, Dependency Track, Grafana, Ingress,~\dots)~entre dans le même jeu que les connecteurs Prometheus~/~Kyverno listés dans le schéma d'environnement validé avec le développement.

\section{Objectif}

Recenser ainsi les~\textbf{cibles d'intégration de l'exécutable}~(CI, registre, GitOps, analyse et dépendances, supervision)~: non pas «~tout savoir du marché~», mais suivre où le développement de~\texttt{paas} passe effectivement lorsque la configuration est disponible.

\subsection{Identifier les outils retenus}

\subsubsection{JENKINS :}

Jenkins est un serveur d’automatisation open-source qui permet d’automatiser diverses tâches liées au développement de logiciels, y compris la construction, le test et le déploiement de code. Il est particulièrement utile dans les environnements de développement qui suivent les principes d’intégration continue et de livraison continue (CI/CD). Les avantages de l’utilisation de Jenkins incluent l’automatisation des tâches de développement, la facilité d’intégration avec d’autres outils, la personnalisation, la scalabilité et le fait qu’il soit open-source. Grâce à Jenkins, les développeurs peuvent améliorer l’efficacité de leur processus de développement logiciel en réduisant les erreurs et en accélérant la livraison de code de qualité supérieure.
\subsubsection{ARTIFACTORY :}

Artifactory est un gestionnaire de dépôts universel qui permet de stocker et de gérer des artefacts logiciels, tels que des fichiers JAR, des images Docker, des packages NPM, des bibliothèques Python, des packages Ruby et bien d’autres encore. Il est utilisé par les équipes DevOps pour stocker, gérer et distribuer des artefacts dans leurs projets et leurs workflows. Cet outil est particulièrement utile, car il simplifie le processus de gestion des artefacts logiciels en fournissant un emplacement centralisé pour stocker et gérer les packages de différents types.

\subsubsection{SONARQUBE :}

SonarQube est un outil open source d’analyse statique de code qui permet de mesurer la qualité du code source. Il aide à détecter les bugs, les vulnérabilités de sécurité, les duplications de code et les violations de conventions de codage. Il prend en charge de nombreux langages de programmation et fournit des rapports détaillés sur la qualité du code, aidant ainsi les développeurs à identifier les problèmes et à améliorer la qualité globale de leurs applications.

\subsubsection{HARBOR :}

Harbor est un registre d’images de conteneurs open source qui permet de stocker, gérer, signer et distribuer des images Docker et OCI. Il fournit une plateforme sécurisée pour gérer les images de conteneurs dans un environnement d’entreprise, avec des fonctionnalités telles que la gestion des accès, la réplication, l’analyse des vulnérabilités et la signature des images. Harbor est conçu pour répondre aux besoins des organisations souhaitant disposer d’une solution fiable et évolutive pour gérer leurs images de conteneurs en toute sécurité.

\subsubsection{JFROG CLI :}

JFrog CLI est un outil en ligne de commande open source développé par JFrog pour gérer les packages et les artefacts sur les serveurs Artifactory. Il permet aux développeurs et aux administrateurs systèmes de télécharger, téléverser, copier, supprimer et rechercher des packages, ainsi que de configurer des pipelines de construction et des workflows de déploiement. Cet outil facilite l’automatisation et la gestion efficace des artefacts logiciels.
\subsubsection{ARGO CD :}

ArgoCD est un outil open source de déploiement continu pour les applications Kubernetes. Il surveille les modifications apportées aux ressources Kubernetes et déploie automatiquement les mises à jour en fonction de la configuration de déploiement spécifiée. ArgoCD permet une livraison continue du code en production, une gestion efficace des versions, une sécurité renforcée, utilise les principes de l’infrastructure as code et est facile à installer et à utiliser.

\subsubsection{DEPENDENCY TRACK :}

Dependency Track est une plateforme open source de gestion des vulnérabilités pour les logiciels. Elle est conçue pour aider les développeurs à identifier, suivre et résoudre les vulnérabilités de sécurité au sein de leurs dépendances logicielles. Dependency Track fournit des informations précieuses sur les vulnérabilités des composants utilisés dans les applications, permettant de prendre des décisions éclairées. En s’appuyant sur plusieurs sources telles que NVD, VulnDB et ExploitDB, elle offre des informations détaillées sur les vulnérabilités connues ainsi qu’un aperçu global de l’état de sécurité des dépendances. Elle permet également d’identifier les composants tiers présentant des risques et de suivre l’évolution de la correction des vulnérabilités.

\subsubsection{KUBERNETES :}

Kubernetes est une plateforme open source de gestion de conteneurs. Elle permet le déploiement, la mise à l’échelle et la gestion d’applications conteneurisées. Largement utilisée dans les architectures de microservices et les environnements cloud, elle simplifie la gestion des conteneurs en proposant une orchestration automatisée : installation, configuration, scaling et monitoring. Kubernetes apporte une solution efficace aux défis liés à la disponibilité, à la gestion de la charge, aux mises à jour et à la mise en réseau .

\subsubsection{GRAFANA :}

Grafana est une plateforme open source d’analyse et de visualisation de données en temps réel. Elle permet de collecter, analyser et afficher des données provenant de différentes sources sous forme de tableaux de bord interactifs. Grâce à ses capacités de personnalisation et à ses nombreux plugins, Grafana offre une grande flexibilité pour surveiller les performances des systèmes et des applications.

\subsubsection{PROMETHEUS :}

Prometheus est une plateforme open source de surveillance et d’alerte pour les systèmes informatiques. Elle collecte et stocke des métriques en temps réel, permettant ainsi d’analyser les performances et de détecter les anomalies. Prometheus utilise un modèle de données basé sur des séries temporelles et s’intègre facilement avec des outils comme Grafana pour la visualisation. Il est largement utilisé dans les environnements Kubernetes pour le monitoring des applications et des infrastructures.

\subsubsection{DEPENDENCY CHECK :}

Dependency Check est un outil open source qui permet d’identifier les vulnérabilités dans les dépendances des applications. Il analyse les fichiers de configuration et compare les bibliothèques utilisées avec des bases de données de vulnérabilités connues telles que NVD. Il génère des rapports détaillés sur les failles détectées et aide les développeurs à corriger les risques de sécurité dans leurs applications. Il supporte différents formats de sortie comme HTML, XML et JSON.
\subsubsection{HAPROXY :}

HAProxy est un outil open source de répartition de charge qui permet de distribuer le trafic réseau entre plusieurs serveurs afin d’assurer la haute disponibilité et la performance des applications. Il utilise les concepts de frontend et backend pour gérer le routage du trafic. Dans un environnement Kubernetes, HAProxy peut être utilisé comme Ingress Controller pour diriger les requêtes vers les services appropriés, en configurant des règles de routage basées sur les ressources Ingress. Il contribue ainsi à une gestion efficace du trafic, à l’amélioration des performances et à la disponibilité des applications.

\subsubsection{NGINX INGRESS CONTROLLER :}

Nginx Ingress Controller est un contrôleur Kubernetes qui permet de gérer les accès externes aux services d’un cluster via l’API Ingress. Il agit comme un point d’entrée unique pour les requêtes HTTP et HTTPS, en les redirigeant vers les services appropriés selon des règles définies. Il est particulièrement adapté aux architectures microservices, car il facilite la gestion du trafic, la répartition de charge, la mise en cache, la compression et la redirection. Il offre une solution flexible et évolutive pour contrôler l’accès aux applications déployées dans Kubernetes.

\subsubsection{OPA Gatekeeper :}

OPA Gatekeeper est un outil basé sur Open Policy Agent (OPA) qui permet de définir et d’appliquer des politiques de sécurité dans Kubernetes. Il utilise le langage déclaratif Rego pour spécifier des règles de conformité et de contrôle d’accès. Grâce à Gatekeeper, il est possible de valider les configurations des ressources Kubernetes et de s’assurer qu’elles respectent les politiques définies, comme les restrictions d’accès ou les bonnes pratiques de sécurité. Il renforce ainsi la sécurité et la gouvernance des clusters Kubernetes et pour metter en place des politiques de conformité reglementaire .
\subsubsection{HELM :}

Helm est un outil open source pour Kubernetes qui facilite la gestion des applications en simplifiant leur déploiement, leur mise à jour et leur configuration. Il utilise des packages appelés \textit{charts}, qui contiennent les ressources nécessaires (fichiers YAML) pour déployer une application. Helm permet ainsi d’assurer une gestion cohérente et reproductible des applications, tout en offrant une grande flexibilité dans la configuration et la distribution. Il est souvent utilisé avec d’autres outils Kubernetes pour améliorer l’automatisation et la gestion des environnements.

\subsubsection{TRIVY :}

Trivy est un outil open source de détection de vulnérabilités pour les conteneurs et les images OCI. Il permet d’identifier rapidement les failles de sécurité présentes dans les images Docker, les dépendances et les configurations. Trivy s’intègre facilement dans les pipelines CI/CD afin d’automatiser l’analyse de sécurité et de détecter les vulnérabilités connues grâce à des bases de données régulièrement mises à jour. Il contribue ainsi à renforcer la sécurité des applications déployées.

\subsubsection{COSIGN :}

Cosign est un outil open source développé par Google qui permet de signer et de vérifier les images de conteneurs afin de garantir leur intégrité et leur authenticité. Il repose sur des mécanismes de cryptographie à clé publique pour sécuriser les images. Cosign permet aux développeurs de vérifier l’origine des images avant leur déploiement et s’intègre facilement dans les workflows DevOps, contribuant ainsi à renforcer la chaîne d’approvisionnement logicielle (software supply chain).
\subsubsection{DOCKER :}

Docker est une plateforme open source qui permet de créer, déployer et exécuter des applications dans des conteneurs. Ces conteneurs regroupent tous les éléments nécessaires à l’exécution d’une application (code, bibliothèques, dépendances), garantissant ainsi une portabilité entre différents environnements. Docker simplifie le processus de déploiement en assurant une isolation des applications et en facilitant leur gestion et leur mise à l’échelle.

\subsubsection{CONTAINERD :}

Containerd est un moteur de conteneurs open source qui gère le cycle de vie des conteneurs, notamment leur création, exécution, arrêt et suppression. Il fournit une interface standardisée pour la gestion des images et des conteneurs, et est utilisé comme runtime par des plateformes comme Kubernetes. Léger et performant, containerd est conçu pour les environnements à grande échelle et offre une grande flexibilité dans la gestion des conteneurs.

\subsubsection{SPRING BOOT :}

Utilisé de façon \textbf{optionnelle} dans le dépôt pour des services Java~(ex.\ proxys d’API vers des outils externes)~; le plan de contrôle principal et les routes API métiers sont portés par \textbf{Next.js~(App Router)~+~Node}.

\subsubsection{NEXT.JS~(FRONT PLAN DE CONTRÔLE)~:}

\textit{Next.js}~(React)~constitue l’IHM~(tableaux de bord, création de projets, suivis)~et les~\emph{handlers} API côté serveur reliant la base~(Prisma)~, la session JWT et les intégrations~(Jenkins, Kubernetes, registre)~.

\subsubsection{TERRAFORM :}

Terraform est un outil open source développé par HashiCorp qui permet de gérer l’infrastructure en tant que code (Infrastructure as Code). Il utilise un langage déclaratif appelé HCL pour définir et provisionner les ressources sur différentes plateformes cloud telles que AWS, Azure ou Google Cloud. Terraform permet de créer des infrastructures reproductibles, de gérer les versions des configurations et d’automatiser les déploiements. Il améliore la collaboration entre les équipes et offre une meilleure visibilité sur les ressources et leurs dépendances.

\subsubsection{GROOVY :}

Groovy est un langage de programmation orienté objet et fonctionnel, souvent utilisé dans l’écosystème Java. Compatible avec la syntaxe Java, il facilite l’intégration avec les bibliothèques existantes. Groovy est couramment utilisé pour écrire des scripts, notamment dans les pipelines CI/CD avec Jenkins. Il permet d’automatiser diverses tâches de développement, de gestion et de déploiement, tout en offrant une grande flexibilité et simplicité d’utilisation.
\subsubsection{AWS :}

Amazon Web Services (AWS) est une plateforme de services cloud proposée par Amazon qui permet de fournir des ressources informatiques à la demande. Elle offre une large gamme de services tels que le calcul, le stockage, les bases de données, le réseau et la sécurité. AWS permet aux entreprises de déployer et gérer des applications de manière flexible, tout en bénéficiant d’une haute disponibilité, d’une scalabilité et d’une fiabilité élevées. Grâce à son modèle de paiement à l’utilisation, AWS constitue une solution adaptée aux besoins variés des projets, allant des startups aux grandes entreprises.

\subsubsection{Calcul du coût AWS :}

Dans le cadre de notre projet, nous avons réalisé une estimation du coût de l’infrastructure AWS en utilisant le calculateur de prix AWS. Cette estimation prend en compte les différentes ressources nécessaires pour héberger et déployer notre solution.

Les composants principaux utilisés sont :

\begin{itemize}
    \item Instances EC2 (Elastic Compute Cloud) avec un système Linux Debian 11
    \item Serveur Jenkins et SonarQube
    \item Serveur Artifactory et Dependency Track
    \item Serveur Harbor
    \item Serveur HAProxy
    \item Cluster Kubernetes (K8s Master et K8s Worker)
\end{itemize}

Afin de garantir la sécurité et l’isolation de l’infrastructure, nous avons configuré un Virtual Private Cloud (VPC) avec des sous-réseaux publics et privés. Des groupes de sécurité ont été définis pour contrôler les accès aux différentes instances.

Cette architecture permet de répondre aux besoins du projet tout en assurant un bon équilibre entre performance, sécurité et coût.
\section{Bilan}

Les fiches précédentes rappellent le rôle de chaque~\textbf{outil pris dans le tableau \texttt{paas/docs/STACK\_OVERVIEW.md}} avant de détailler l'infra de démo sous AWS~(coût indicatif)~: elles~\emph{instrumentent}~\texttt{paas/frontend}. La sélection n'est pertinent que lorsqu'une URL ou un jeton est injecté depuis~\texttt{.env} et consommé par les services~\texttt{src/server/*/}.
